% Generated draft: requires mathematical and source review.
\documentclass[11pt]{article}
\usepackage[margin=25mm]{geometry}
\usepackage{fontspec}
\setmainfont{Latin Modern Roman}
\usepackage{amsmath,amssymb,mathtools,unicode-math}
\setmathfont{Latin Modern Math}
\usepackage{enumitem,needspace,xurl,hyperref}
\setlength{\parindent}{0pt}
\setlength{\parskip}{5pt}
\setlength{\emergencystretch}{3em}
\newcommand{\N}{\mathbb{N}}
\newcommand{\Z}{\mathbb{Z}}
\newcommand{\Q}{\mathbb{Q}}
\newcommand{\R}{\mathbb{R}}
\newcommand{\C}{\mathbb{C}}
\newcommand{\F}{\mathbb{F}}
\newcommand{\A}{\mathbb{A}}
\newcommand{\PP}{\mathbb{P}}
\newcommand{\E}{\mathbb{E}}
\newcommand{\eps}{\varepsilon}
\newcommand{\dd}{\,\mathrm{d}}
\begin{document}
\section*{UnsolvedMath ID: 20001913}
\textbf{Source problem code:} catalogue entry
\subsection{Definitions and mathematical statement}
Let \(M\) be smooth, closed and connected, with \(n=\dim M\geq3\), as specified by the source's surrounding conventions. Let \(\mathcal G\) be its intended space of Riemannian metrics. The provided verification breaks off before defining the topology on \(\mathcal G\), so density must retain that unresolved choice. A geodesic flow is transitive in the source's sense if some unit geodesic has tangent vectors dense in the whole unit tangent bundle:
\[
 \exists v\in S_gM:\quad
 \overline{\{\phi_g^t(v):t\in\mathbb R\}}=S_gM.
\]
Question 8.4.2 asks whether, in the intended topology, every metric neighborhood contains a metric with this property.

The source contrasts dimension two: suitable integrable metrics have irrational Liouville tori persisting under sufficiently small perturbations by KAM theory, producing nontransitive neighborhoods. The exact nondegeneracy, arithmetic and regularity hypotheses behind this statement are omitted. It also relates the higher-dimensional density question to Arnold diffusion. These observations do not specify a higher-dimensional solution.

The paired aid records two necessary diagnostics. First, a continuous first integral \(I:S_gM\to\mathbb R\), satisfying \(I\circ\phi_g^t=I\), must be constant if a dense orbit exists. In particular, let \(K\) be a symmetric covariant Killing tensor of degree \(m\geq1\), so \(K(v,\ldots,v)\) is constant along geodesics. Transitivity forces
\[
 K=0\quad(m\text{ odd}),\qquad
 K=c\,g^{\odot r}\quad(m=2r),
\]
where \(c\) is a global constant and \(\odot\) denotes symmetric tensor product. Absence of nontrivial such tensors is necessary, not sufficient, for transitivity.

Second, in an integrable \(n\)-degree-of-freedom geodesic system, regular Liouville tori have dimension \(n\), while the unit-energy shell has dimension \(2n-1\). Their codimension is \(n-1\). The aid reports that the complement of a finite family of embedded regular Liouville tori in the connected unit shell is path connected for \(n\geq3\), since their codimension is at least two. On a regular open part the same statement is understood componentwise in a connected ambient part. In dimension two the codimension is one, so the finite separating-barrier mechanism is possible there.

This proves only that the finite codimension-one barrier argument does not transfer verbatim to higher dimensions. It does not control an infinite Cantor family of surviving KAM tori, destroy exact first integrals, or produce an orbit moving through resonance channels. The supplied literature report leaves the original density and associated diffusion problems open.
\subsection{Short English statement}
The source asks whether transitive geodesic flows are dense among metrics in dimension at least three. The aid supplies necessary conditions and a finite-barrier distinction: a dense orbit forbids nonconstant continuous integrals, while finitely many higher-codimension Liouville tori cannot separate the energy shell. Neither fact establishes dynamical passage or the requested density.
\subsection{Sources}
\begin{itemize}
\item Local source record: catalogue entry, catalogue entry; canonical \texttt{aim-geometry-notes.json}, \texttt{notes[250]}, source item 8.4, extracted Question 8.4.2. Title: ``A first-integral obstruction and finite Liouville-barrier dichotomy.''
\item Keith Burns and Vladimir S. Matveev, \emph{Open problems and questions about geodesics}, AIM workshop ``Geodesics,'' supplied reference \url{https://aimath.org/pastworkshops/geodesicsproblems.pdf}. This draft uses the supplied text without accessing a PDF.
\item Original context invokes KAM theory for the surface contrast and Arnold diffusion for the higher-dimensional question, without full bibliographic entries. The first-integral/Killing-tensor and finite-torus diagnostics are candidate contributions of research attempt 1.
\item Registry dated 14 August 2026: ``partially\_solved,'' ``PARTIAL-PROGRESS,'' statement status ``exact''; attempt partial, research-run status open, rigor ``proved,'' novelty confidence low. The verification restores smooth closed connected conventions but is truncated before the metric-space definition. The machine-generated, non-peer-reviewed aid requires independent verification before citation.
\end{itemize}
\end{document}
